% ================= 宏包 =================

\usepackage{fontspec}
\setmainfont{Times New Roman}           % 英文字体
\setCJKmainfont{SimSun}[BoldFont=SimHei]            % 中文字体：宋体



\usepackage{amsmath,amssymb,amsthm}
\usepackage{mathrsfs}   % 花体字母
\usepackage{tikz-cd}




\usepackage[top=2.54cm, bottom=2.54cm, left=3.17cm, right=3.17cm, bindingoffset=1.0cm]{geometry}




% 行距：1.5倍
\usepackage{setspace}
\onehalfspacing
% 页眉页脚
\usepackage{fancyhdr}
\setlength{\headheight}{13pt}
\pagestyle{fancy}
\fancyhf{}                              % 清空原有设置
\renewcommand{\headrulewidth}{0pt}
\fancyfoot[C]{\zihao{5}\thepage}
\fancypagestyle{plain}{
    \fancyhf{}
    \renewcommand{\headrulewidth}{0pt}
    \fancyfoot[C]{\zihao{5}\thepage}
}



\ctexset{
    chapter = {
      format = \zihao{3}\bfseries\raggedright,
      name = {,},
      number = \arabic{chapter},
      beforeskip = 0pt,
      afterskip = 20pt,
     },
    section = {
      format = \zihao{4}\bfseries\raggedright,
      name = {,},
      number = \arabic{chapter}.\arabic{section},
      beforeskip = 0pt,
     },
    subsection = {
            format = \zihao{-4}\bfseries\raggedright,   % 小四号
            name = {,},
            number = \arabic{chapter}.\arabic{section}.\arabic{subsection},
        }
}




\usepackage[titles]{tocloft}
\usepackage{setspace}

% 设置目录标题格式（三号宋体加粗居中）
\renewcommand{\cfttoctitlefont}{\hfill\zihao{3}\bfseries\heiti}  
\renewcommand{\cftaftertoctitle}{\hfill}
\renewcommand{\cftbeforetoctitleskip}{0pt}
\renewcommand{\cftaftertoctitleskip}{20pt}   % 标题与内容间距


% 设置目录项字体为小四号
\renewcommand{\cftchapfont}{\zihao{-4}\normalfont}      % 一级标题（章）
\renewcommand{\cftsecfont}{\zihao{-4}\normalfont}       % 二级标题（节）
\renewcommand{\cftsubsecfont}{\zihao{-4}\normalfont}    % 三级标题（小节）


% 设置页码与标题之间的点线
\renewcommand{\cftchapleader}{\cftdotfill{\cftdotsep}}   % 一级标题点线
\renewcommand{\cftsecleader}{\cftdotfill{\cftdotsep}}
\renewcommand{\cftsubsecleader}{\cftdotfill{\cftdotsep}}



\usepackage{hyperref}
\usepackage{enumitem}
\usepackage[backend=biber,style=numeric,sorting=none]{biblatex}
\addbibresource{refs.bib}



% ================= 定理环境 =================
\theoremstyle{plain}
\newtheorem{theorem}{定理}[section]
\newtheorem{proposition}[theorem]{命题}
\newtheorem{lemma}[theorem]{引理}
\newtheorem{corollary}[theorem]{推论}
\newtheorem{example}[theorem]{例}
\newtheorem{conjecture}[theorem]{猜想}

\theoremstyle{definition}
\newtheorem{definition}[theorem]{定义}

\theoremstyle{remark}
\newtheorem{remark}[theorem]{注}
% ========== 自定义数学命令 ==========
\newcommand{\Int}{\operatorname{Int}}
\newcommand{\loc}{\mathrm{loc}}
\renewcommand{\d}{\, \mathrm{d}}
\newcommand{\supp}{\operatorname{supp}}

\newcommand{\rank}{\operatorname{rank}}

\renewcommand{\Im}{\operatorname{Im}}
\newcommand{\id}{\operatorname{id}}

\newcommand{\Ker}{\operatorname{Ker}}
\newcommand{\Coker}{\operatorname{Coker}}
\newcommand{\Coim}{\operatorname{Coim}}
\newcommand{\Codim}{\operatorname{Codim}}
\newcommand{\Tor}{\operatorname{Tor}}
\newcommand{\Ann}{\operatorname{ann}}
\newcommand{\Gal}{\operatorname{Gal}}
\newcommand{\ord}{\operatorname{ord}}

\newcommand{\Obj}{\operatorname{Obj}}
\newcommand{\Mor}{\operatorname{Mor}}

\newcommand{\Jac}{\operatorname{Jac}}
\newcommand{\Rad}{\operatorname{Rad}}
\newcommand{\Nil}{\operatorname{Nil}}
\newcommand{\lcm}{\operatorname{lcm}}

\newcommand{\Hom}{\operatorname{Hom}}
\newcommand{\End}{\operatorname{End}}
\newcommand{\Aut}{\operatorname{Aut}}
\renewcommand{\det}{\operatorname{det}}

\newcommand{\GL}{\operatorname{GL}}



\newcommand{\ind}{\operatorname{ind}}


\newcommand{\gr}{\operatorname{gr}}

\newcommand{\Set}{\mathbf{Set}}
\newcommand{\Top}{\mathbf{Top}}
\newcommand{\op}{\mathrm{op}}
\newcommand{\Grp}{\mathbf{Grp}}
\newcommand{\Ab}{\mathbf{Ab}}
\newcommand{\Rng}{\mathbf{Rng}}
\newcommand{\CRing}{\mathbf{CRing}}
\newcommand{\Mod}{\mathbf{Mod}}
\newcommand{\Ring}{\mathbf{Ring}}
\newcommand{\Vect}{\mathbf{Vect}}
\newcommand{\Sh}{\mathbf{Sh}}
\newcommand{\Open}{\mathbf{Open}}

\newcommand{\Ext}{\operatorname{Ext}}

\newcommand{\res}{\operatorname{res}}
\newcommand{\dist}{\operatorname{dist}}
\newcommand{\disc}{\operatorname{disc}}
\newcommand{\Nat}{\operatorname{Nat}}
\newcommand{\Nm}{\operatorname{Nm}}

\newcommand{\Frac}{\operatorname{Frac}}
\newcommand{\Spec}{\operatorname{Spec}}
\newcommand{\Der}{\operatorname{Der}}
